% 机械波及叠加（高中）
% license Usr
% type Tutor
\pentry{机械振动（高中）\nref{nod_HSPM09}}{nod_6f22}
\subsection{机械波}
机械波是机械振动在空间中传播的现象。称振动的源头为波源，称承载振动的物质为介质\footnote{真空实际并非空无一物，而是能量最低的一个状态。实际上，从量子物理的观点来看，无论真空能是否为0（我们可以灵活地调整能量零点），真空总是会在极短时间内产生和湮灭粒子-反粒子对。这种涨落会被时间“抹平”，所以在相当长一段时间内观察，能量平均值还是基态。但极短时间的这种量子涨落亦能传播振动，且已经为实验所证实——\href{https://www.rmzxw.com.cn/c/2023-08-25/3399291.shtml}{真空传声}}。从微观上看，只要有波源在作振动，且质点之间有弹性恢复力，那么必然引起传播方向上，下一个相邻质点作振动。每一个质点，都可以看作接下来的，机械振动的波源，如果没有其他能量损耗，机械波的传播是没有止境的。其传播路径上的每个质点都可视为弹簧振子，在受到上一个质点的扰动后，才开始在平衡位置附近作振动。

从上述过程来看，机械波对振动能量的传播不言而喻。实际上，机械波还是振动状态（相位）的传播。假设振动在传播过程中没有损耗，如果A点为较近点，B点是较远点，且振动从A传播到B需要时间$t$，那么B在某时刻$t_0$的相位总是B在$t_0-t$的相位。所以不同质点的振动图像都有共同属性如周期，振幅，有所不同的只是相位。
\subsubsection{简谐波的描述与图像}
如果我们把机械波简化成一连串质点的运动，且设波源作简谐振动，介质是理想介质（即波速稳定），那么在任意时刻，给这一串质点“拍照”，它们的位置构成简谐函数的波形。
比如设波源的振动形式为$y=\sin(\omega t)$，波速为$v$，那么位于波源正前方$x$的质点位移为
\begin{equation}\label{eq_mcwave_1}
y=\sin [\omega (t-\frac{x}{v})]~.
\end{equation}
显然如果在某时刻拍照，固定$t$后的质点位移$y$是关于$x$的简谐函数，且周期与振动周期相同，常用$T$表示。%%插图，振动和波动

我们称波在一个周期传播的距离为\textbf{波长}，常用$\lambda $表示，满足$\lambda=vT$。波源的振动决定周期，或者说频率，所以无论机械波经过何种介质，其周期都是不变的。然而，介质的物理性质决定了波速，所以改变介质就会改变波速和波长。
对应山峰和山谷，我们称波的位移最大值为\textbf{波峰}，最小值为\textbf{波谷}。对于简谐波，波峰和波谷的相位差总是$\pi$。

\subsection{波的叠加}
在普通物理学里，我们通常假设介质是理想线性的。这意味着，一方面，机械波在传播过程中没有能量损耗，另一方面，机械波的传播不会改变质点的振动响应，其结果是，传播路径上的每个质点都等效于一个弹性系数相同，振幅不变的弹簧，这就有了\autoref{eq_mcwave_1} 的振动方程。此外，如果有两列波在空间传播，我们会额外假设传播是线性的，也就是不同波在空间中的叠加不会影响彼此，更不会影响介质总的响应。这就好比用不同的力拉一根弹簧，假设它不会崩坏。

如果不同简谐波都满足以下条件：
\begin{enumerate}
\item 频率相同
\item 相位差恒定\footnote{如果波源运动，相位差就可能不恒定}
\item 有沿着同一方向的振动分量
\end{enumerate}
此时，我们就称波的叠加为\textbf{波的干涉}。其特点是，任意一处的质点都在做简谐振动。下面我们来证明这个结论。



在机械振动一节中，我们通过推导，证明了匀速圆周运动的位移矢量投影到坐标轴后，正是简谐振动，所以简谐振动和匀速圆周运动是一一对应的，简谐振动的角频率是匀速圆周运动的角速度，振幅则是匀速圆周运动的半径。现在假设两个波源传播到某一位置后，在某一方向的振动方程分别为：
\begin{equation}\label{eq_mcwave_2}
\begin{aligned}
y_A&=A\sin(\omega t+\phi_1)\\
y_B&=B\sin (\omega t+\phi_2)~,
\end{aligned}
\end{equation}
如下图所示。由正交分解法可知，合振动位移正是匀速圆周运动两个位移矢量合成后在相同坐标轴下的投影。于是我们立刻能发现，由于两个振动的相位差恒定，合振动作振幅不变的简谐振动。

\begin{figure}[ht]
\centering
\includegraphics[width=6cm]{./figures/c5f27b2b59b3e55e.png}
\caption{角频率相同，相位差恒定的简谐振动合成} \label{fig_mcwave_1}
\end{figure}
那么自然而然的，合振动的振幅取决于两个位移矢量合成后的模长，当两个位移矢量同向，也就是波源在该位置的相位差为$2n\pi(n\in \mathbb{N}_0)$时，合振动振幅最大，我们把质点所在位置称之为振动加强点；当位移矢量反向，即相位差为$(2n+1)\pi$时，合振动振幅最小，此时位置则是振动减弱点。读者可验证，该结论与两个波源在此处的振动方向无关。即使一个振动方向先向上，后向下（相当于质点逆时针转动），另一个振动方向先向下，后向上（相当于质点顺时针转动），合振幅规律与振动方向相同的别无二致。
\subsubsection{另一证明过程}
展开\autoref{eq_mcwave_2} ，我们有
\begin{equation}
\begin{aligned}
& y_A=A\left(\sin \omega t \cos \phi_1+\cos \omega t \sin \phi_1\right), \\
& y_B=B\left(\sin \omega t \cos \phi_2+\cos \omega t \sin \phi_2\right)~.
\end{aligned}
\end{equation}
，展开后有

\begin{equation}
y=\left(A \cos \phi_1+B \cos \phi_2\right) \sin \omega t+\left(A \sin \phi_1+B \sin \phi_2\right) \cos \omega t\equiv C\sin (\omega t+\phi)~,
\end{equation}
则合振动振幅满足

\begin{equation}
\begin{aligned}
C^2
= & \left(A \cos \phi_1+B \cos \phi_2\right)^2+\left(A \sin \phi_1+B \sin \phi_2\right)^2 \\
= & A^2\left(\cos ^2 \phi_1+\sin ^2 \phi_1\right)+B^2\left(\cos ^2 \phi_2+\sin ^2 \phi_2\right) \\
& +2 A B\left(\cos \phi_1 \cos \phi_2+\sin \phi_1 \sin \phi_2\right) \\
= & A^2+B^2+2 A B \cos \left(\phi_1-\phi_2\right)~,
\end{aligned}
\end{equation}
上述结论得证。
\subsection{相位差与波程差}
假设两个波源相干，且波长相同，$P$点距波源$A,B$的距离分别是$r_1,r_2$，那么两个波源在$p$处引起的振动方程分别是
\begin{equation}
\begin{aligned}
y_A&=\sin[\omega t-kr_1+\phi_1]\\
y_B&=\sin[\omega t-kr_2+\phi_2]~,
\end{aligned}
\end{equation}
其中$k=2\pi/\lambda$。则相位差为
\begin{equation}
\Delta\phi=k(r_1-r_2)+\phi_2-\phi_1~.
\end{equation}

中学物理常引入“振动步调”的概念来描述波源的振动相位。振动步调即振动位移，那么“振动步调一致”即波源的振动位移方程相同，$\phi_2-\phi_1=0$。当相位差为$2n\pi$，也就是振动加强时，波程差$r_1-r_2=n\lambda$，当振动减弱，也就是相位差为$(2n+1)\pi$时，有$r_1-r_2=(2n+1)\lambda/2$。

若两个波源振动步调相反，则$\phi_2-\phi_1=\pi$。同理可得，振动加强时$r_1-r_2=(2n-1)\lambda/2$，振动减弱时$r_1-r_2=2n\pi$，观察可知，条件与步调一致时相反。


